
\documentclass[12pt]{article}
\usepackage[letterpaper,margin=1.0in]{geometry}
\usepackage[T1]{fontenc}
\usepackage[utf8]{inputenc}
\usepackage{lmodern}
\usepackage{microtype}
\usepackage{setspace}
\usepackage{parskip}
\usepackage{needspace}
\usepackage[hidelinks]{hyperref}
\setstretch{1.12}
\begin{document}
\begin{titlepage}
\centering
\vspace*{1.5in}
{\Huge\bfseries How Order Holds\par}
\vspace{0.4in}
{\Large On Active Boundaries, Self-Maintaining Systems, and the Cost of Staying Organized\par}
\vspace{0.8in}
{\large Andrew Poinçot\par}
{\large Cripsis Logic\par}
\vfill
{An essay introducing the academic paper that follows.\par}
{Revised draft v10.9.2: May 8, 2026 \par}
{Public edition.\par}

\end{titlepage}


Why does anything stay together? Not metaphorically --- literally. A living cell is mostly water, salt, and a tangle of long-chain molecules that have no particular reason to remain organized. Left to its own devices, the universe pulls things apart. The second law of thermodynamics is uncompromising about this: in any closed system, disorder increases. And yet, for a few decades at a time, that bag of salty water stays organized enough to read this sentence. So does a redwood. So does a coral reef, a city, an economy, the magnetic field of a star, the routing system that delivered this essay to your screen. How? The honest answer, the one physicists have been circling since the nineteenth century, is that none of these things are closed systems. They each maintain their structure by doing active work --- taking in energy or information from outside, exporting disorder back to the environment, and burning the difference to hold their shape. A cell metabolizes. A redwood photosynthesizes. A city imports food and exports trash. A star fuses hydrogen against gravitational collapse. A computer system handling a million requests per second runs algorithms that decide what to admit, what to defer, and what to reject, all in service of not falling over. These examples sound like very different problems. They are. The machinery a cell uses to stay alive bears little resemblance to the machinery a star uses to stay lit, which bears little resemblance to the machinery a software system uses to stay responsive. The microscopic mechanisms diverge wildly. But at a higher level of description --- the level a physicist would call phenomenological, meaning concerned with patterns of behavior rather than with underlying mechanism --- these systems share something striking. They all exhibit thresholds. They all have regimes where they hold their shape and regimes where they don't. They all, when pushed too far, collapse suddenly rather than gradually. And the line between holding and collapsing tends to be sharper than intuition would suggest. This essay is about that line, and about a small contribution to understanding where it lives. It will not assume any background in physics or mathematics. It will assume a reader who finds the question interesting and is willing to spend an hour thinking it through.


\section*{I. The Long Conversation}

The question of how order persists in systems far from equilibrium has a long pedigree. It is, in some sense, the central question of twentieth-century physics outside the relativity-and-quantum mainstream --- the question that occupied the people who looked at the universe and asked not what it is made of, but how it holds itself up. Erwin Schr{\"o}dinger, the physicist who gave us the wave equation that underwrites quantum mechanics, spent the last years of his life writing a small book called What Is Life? In it, he argued that living things stay alive by feeding on what he called negative entropy --- by extracting order from their environment to compensate for the disorder they generate internally. The book is brief, speculative, and occasionally wrong, but it set the agenda for a generation of biophysicists and gave us the conceptual vocabulary we still use. Schr{\"o}dinger's central insight was that life is not a violation of thermodynamics but an expression of it: living systems are local regions of decreasing entropy, paid for by larger increases in entropy elsewhere. The cost is real. The accounting balances. Ilya Prigogine, working in Brussels a generation later, formalized this into the mathematics of dissipative structures. Prigogine showed that when energy flows through a system far enough from equilibrium, ordered patterns can emerge spontaneously --- convection cells in a heated fluid, chemical clocks that oscillate between states, the Belousov--Zhabotinsky reaction with its hypnotic spiral waves. Prigogine won the Nobel Prize in 1977 for this work. His central message was that order is not the opposite of chaos; under the right conditions, order is what chaos produces. Stuart Kauffman, at the Santa Fe Institute, took this further. In a series of provocative books beginning in the late 1980s, Kauffman argued that life sits at what he called the edge of chaos --- a regime where systems are organized enough to function but disordered enough to adapt. Too much order produces frozen, brittle systems unable to respond to change. Too much chaos produces noise. The interesting regime, Kauffman claimed, is the narrow band between, and evolution drives systems toward it. Whether this is literally true remains debated. That it is suggestive remains undeniable. More recently, Karl Friston, a British neuroscientist, has framed perception and cognition themselves as a special case of this larger question. Brains, in Friston's account, stay organized by minimizing the difference between what they predict about the world and what they observe. This minimization is itself an active process --- the brain feeds on information the way cells feed on glucose, and the suppression of surprise is metabolically expensive. Friston's free-energy principle is controversial, dense, and possibly the most ambitious attempt yet to unify thermodynamics, information theory, and biology under a single mathematical roof. Around these named figures has grown a broader tradition --- sometimes called complexity science, sometimes complex adaptive systems, sometimes simply the Santa Fe school --- that has pursued the question across biology, economics, ecology, and computation. Researchers in this tradition have produced a generation of insights about emergence, self-organization, adaptive complexity, networks, and the strange regularities that appear when many simple agents interact under constraint. Three more specific research programs deserve mention because the paper that follows engages with them directly. Friston and his collaborators have developed the formalism of Markov blankets --- statistical boundaries that separate a system's internal states from its environment --- as the mathematical object on which active inference operates. The viability theory tradition, beginning with Jean-Pierre Aubin's 1991 monograph, formalized self-maintenance through the geometry of viability kernels in control systems. And in systems biology, James Sethna, Ilya Nemenman, and their collaborators have shown that complex adaptive networks often reduce, through manifold-boundary approximations of their microscopic kinetics, to a small number of dimensionless ratios --- sometimes just one. Each of these programs answers a version of our central question with different machinery. Each is a serious body of work. None of them quite produces the kind of compact, falsifiable phenomenological law this essay is leading toward, though all of them frame what such a law would have to do. Each of these framings captured part of the phenomenon. Each of them fell short of the same thing: a phenomenological law compact enough to be tested against measurement in a specific system, then extended only as the data warranted. The frameworks were rich. The vocabulary was evocative. The predictions were rare. This is not a complaint, exactly. The work was hard, the systems were varied, and the temptation to generalize was strong. But popular science writing about complex systems has long had a problem: it tends to substitute metaphor for prediction. A theory that explains why an ant colony, a financial market, and a human brain are all examples of the same deep principle is exciting to read about and difficult to test. A theory that says here is the specific number you should measure, and here is what you should see if I am right, and here is what you should see if I am wrong is much less exciting and much more useful. What follows is an attempt at the second kind.


\section*{II. A Question Worth Asking Precisely}

Here is the question, sharpened. Take any system that does active work to maintain its organization --- a cell, a neural circuit, a software service, a market-making operation, a small team trying to ship a product. Suppose we leave that system alone. We don't add new energy, we don't add new information, we don't intervene from outside. Will the system sustain its organized state, or will it collapse to something inert? Phrased this way, the question admits a sharp answer for a surprisingly broad class of systems. The answer is that there exists, for each system, a single dimensionless number --- a ratio of how strongly the system feeds back on itself versus how strongly it dissipates internal complexity --- that determines viability. Below a critical value, the system cannot hold its organization. Above the critical value, it can. The transition between these regimes is not gradual. It is a threshold. This may sound similar to a number readers will have heard about during the recent pandemic: the basic reproduction number, R-zero, which determines whether an infectious disease will spread or die out. The structural similarity is real. Both numbers are dimensionless ratios that cross one as the qualitative behavior of the system changes. Both correspond to thresholds in the underlying mathematics. But the mechanism is entirely different. R-zero is about transmission between hosts. The number that governs the question we're asking here --- call it R, plain --- is about a single system's ability to feed back on itself faster than it dissipates the complexity its own activity produces. The two numbers share a logical shape, not a physical origin. Why should such a number exist? Because the mathematics of self-maintenance, when written carefully, has a specific structure. Imagine a system with two coupled quantities. The first is how much organized, useful activity the system has at any given moment --- call this the activity. The second is how much disorder, complexity, or internal pressure has built up --- call this the entropy, in a deliberately broad sense. These two quantities feed each other. Activity generates internal pressure, because every action produces consequences that have to be processed. Pressure drives further activity, because organized work is often a response to internal complexity that demands attention. But both quantities also decay over time if nothing replenishes them. Activity dissipates without ongoing drive. Pressure relaxes if there's nothing generating new of it. The question of whether the system can sustain itself is the question of whether this two-way coupling is strong enough to overcome the dissipation. If the feedback is weak relative to the decay, the system winds down to nothing. If the feedback is strong relative to the decay, the system can hold a stable, bounded level of activity indefinitely. The viability number R captures this comparison in a single quantity. If R is less than one, dissipation wins. If R is greater than one, feedback wins, and a stable active state becomes available. There is one more piece of mathematical machinery the model needs. Without it, the equations would predict that activity grows without bound --- runaway feedback, infinite escalation, the kind of behavior that real systems do not exhibit because real systems have limits. Cells run out of substrate. Software systems hit hard rate caps. Stars exhaust their fuel. To capture this, the model includes a saturation term: a nonlinear effect that becomes important only when activity gets large, and that bends the trajectory back from runaway. The specific form of this saturation determines exactly how the active state scales with how far above threshold the system is, but the qualitative story is the same regardless: above threshold, the system settles into a stable, bounded active state, and the size of that state grows in a specific predicted way as you push the system further from the threshold.


\section*{III. The Shape of the Answer}

When this is all written down carefully, three predictions fall out, and they are predictions you could in principle test against any real system you suspect the framework applies to. The first prediction is the threshold itself. Sweep R across one --- by tuning whatever parameter in the system controls the feedback --- and the system should transition sharply from quiescent to active. Below the threshold: no spontaneous activity. Above it: bounded, stable activity. The transition should be visible without elaborate statistical machinery, simply by plotting the system's organized activity against the value of R. The second prediction is the amplitude scaling. Above the threshold, the size of the active state should grow as the square root of how far you are above one. Equivalently, the square of the activity should grow linearly with distance above threshold. This is a quantitative claim. It is testable. A system that crosses the threshold but doesn't show this specific scaling has falsified the model --- or at least the specific saturation form the model assumes. The third prediction is the most distinctive, and the one that connects this work to a phenomenon readers may have encountered in popular science before. It is called critical slowing, and it works like this. Near the threshold, when the system is just barely able to maintain itself, its recovery from disturbance becomes arbitrarily slow. Push the system a little --- small perturbation, nothing dramatic --- and watch how long it takes to return to its stable state. Far above threshold, it returns quickly. Just above threshold, it returns slowly. Right at threshold, it doesn't return at all in any practical sense; the recovery time diverges to infinity. This is not a curiosity. Critical slowing has been observed in ecosystems heading toward collapse, in lakes transitioning between clear and turbid states, in financial markets in the months before crashes, and in human brains transitioning between states of consciousness or between healthy and seizure activity. It is one of the more reliable signatures we have that a system is approaching a qualitative transition. The framework here predicts it from first principles, with a specific functional form: the recovery time should grow as one over the distance above threshold. This is a sharp, checkable claim. These three predictions --- the threshold, the amplitude scaling, the critical slowing --- are what the framework commits to. If a calibrated real system exhibits all three with the right functional forms, the framework fits that system. If it exhibits none of them, the framework is wrong. If it exhibits some but not others, the failure modes tell us specifically what kind of extension the framework needs --- different saturation, additional state variables, delayed feedback, and so on. The framework is, in other words, falsifiable in a precise way, and informative even when it fails.


\section*{IV. The Discipline of Being Wrong}

It is worth dwelling on this last point, because it is the part of the work I think matters most, and it is the part most likely to be underappreciated. The history of cross-disciplinary phenomenological theory is littered with frameworks that explained too much. They were elegant, evocative, often profound --- and they accommodated almost any observation, which is to say, they predicted almost nothing. The philosopher Karl Popper argued, famously, that a theory which cannot in principle be falsified is not a scientific theory at all. He was thinking primarily of psychoanalysis and Marxism, but the warning applies more broadly. Theories of complexity, in particular, have a tendency to drift toward the unfalsifiable. They explain why systems are complex by appealing to complexity. They predict that systems will self-organize, and when systems do something else, that too is self-organization. The way to do better is to commit, in advance, to specific things you would not see if you were right, and to admit clearly when you see them. The framework presented here does this. It includes an explicit table of failure modes --- observations that, if seen in a calibrated system, would either falsify the model or tell us exactly what kind of extension is needed. Sustained activity below threshold? Falsified, unless we missed external forcing or a hidden state variable. Persistent oscillations around the active state? The minimal model is incomplete and needs delayed feedback or an additional state. No critical slowing near threshold? Either we mis-estimated R, or the transition isn't the one the model describes. This kind of explicit failure-mode mapping is unusual in popular treatments of complexity science, and it is the part of the paper I would point a skeptical reader to first. The mathematics is well-established, drawn from a long tradition of nonlinear dynamics. The interpretive contribution is genuine but modest. The discipline of being clear about what would prove the framework wrong is, I think, the part that earns the rest of it. A framework like this becomes scientifically useful only if it can be wrong. The framework presented here can be wrong in six specific ways. Each way of being wrong tells us something. That is the whole structure.


\section*{V. Where This Touches the World}

The most immediate application is in software systems that regulate their own behavior under load. This sounds dry, but it shouldn't. As we deploy artificial intelligence at scales that strain the systems carrying it, the question of whether our control infrastructure is doing real regulatory work or merely adding overhead becomes economically and operationally consequential. Concretely: when you send a request to a large AI model, that request doesn't go directly to the model. It passes through a series of regulatory layers --- admission controllers, routing logic, backpressure systems --- that decide whether to serve your request immediately, queue it, defer it, or reject it. These regulatory layers are themselves active boundaries, in the precise sense the framework describes. They take in workload, generate internal complexity in the form of queue pressure and decision uncertainty, and dissipate that complexity through completion, retries, and cleanup. They have all the structural features of the systems the framework was designed to describe. They also produce, as a routine matter of operation, the telemetry needed to test the framework's predictions. Queue depth, decision rates, recovery times --- these are logged in nearly every production system at scale. The framework gives operators of such systems a quantitative answer to a question they currently answer with intuition: is our control layer actually doing useful work? Is it stabilizing the system, or is it merely adding overhead without protective value? Given the cost of running large AI models --- measured in millions of dollars per month for any serious deployment --- the difference between a control layer that genuinely regulates and one that merely consumes is substantial. Beyond this immediate application, the framework plausibly applies wherever the question "is this system maintaining itself, or is it being maintained?" is meaningful. Neural circuits operate near criticality, and the framework's predictions about critical slowing and amplitude scaling have natural neural analogs. Ecological systems exhibit threshold transitions and recovery-time signatures that the framework predicts. Organizations and institutions undergo discontinuous transitions in operational coherence that look, structurally, like the active-to-quiescent transitions the model describes. Markets, particularly market microstructure and clearing infrastructure, regulate organized trading activity against pressure from order flow. Each of these would require its own calibration work, its own choice of proxies for activity and complexity, its own validation. The paper is explicit that these are research directions, not completed claims. The temptation to gesture broadly at universal applicability has been the downfall of many complexity frameworks, and this one tries to avoid the trap by naming specifically where the work has been done --- synthetic demonstrations of the predicted signatures in software-style telemetry --- and where it hasn't. There is also a longer-range direction worth flagging, though it lies outside the present work and is mentioned only as a structural resonance. The active-boundary description --- a boundary that retains information while exchanging entropy with its environment under regulatory constraint --- has a formal echo in black-hole horizon physics. Black-hole horizons are themselves active boundaries in this sense: they retain information about what fell in, exchange entropy with their cosmic environment via Hawking radiation, and obey laws that look thermodynamic. Whether the framework here has anything to say about gravitational systems is a question for a different paper and a different set of derivations. But the structural similarity is real enough to flag, and curious enough to be worth flagging.


\Needspace{9\baselineskip}
\section*{VI. What the Work Doesn't Do}

It is worth being explicit about the limits of what the framework claims, because those limits are part of why the work is offered in the spirit it is. It is not a microscopic theory. The framework does not derive its equations from underlying mechanisms in any specific domain. It is phenomenological --- a description of patterns of behavior --- and its parameters must be measured or estimated for each system to which it is applied. The viability number R is a ratio of quantities that have specific meanings in any given system, but those meanings have to be supplied by the practitioner, not by the framework. It is not a universal theory of complexity. It applies to systems where the active-boundary description fits. There are surely systems for which it does not fit --- systems with strong oscillatory dynamics the minimal model cannot produce, systems with multiple stable states, systems where the relevant nonlinearity is not the cubic saturation the framework assumes. The paper is explicit about these cases and provides the failure modes that would diagnose them. It is not a substitute for empirical work in any specific domain. Calibrating the framework to a real cell, a real ecosystem, a real neural circuit, or a real organization is a research program, not a consequence of having read this paper. The synthetic demonstrations in the paper show what the predicted signatures look like in telemetry-shaped data, but synthetic demonstrations are not real calibrations and the paper does not claim otherwise. And it is not, despite its structural resonances with much grander questions, a theory of life or consciousness or markets or civilization. It is a small mathematical claim about a class of systems that exhibit a specific kind of self-maintenance, with a specific viability threshold and a specific set of falsifiable signatures. It is offered as one node in a long conversation, not as the conversation's conclusion.


\section*{VII. The Invitation}

Why publish a paper like this? Why write an essay introducing it to readers who are unlikely to read the technical version? The honest answer is that work like this has value only if it circulates. A phenomenological framework gains traction when practitioners try it on their own systems. A model becomes useful when other people calibrate it, extend it, falsify it, or replace it with something better. None of that happens if the paper sits on a preprint server unread. More than that, the question the framework addresses is one that matters. We are entering a period where the systems we build are increasingly active boundaries in their own right. Large AI models regulate their own load. Distributed systems negotiate their own state. Autonomous agents make their own decisions about when to act and when to wait. The discipline of asking, of any such system, whether it can sustain itself and at what cost, is going to become more important rather than less. The vocabulary for asking that question precisely has been developing for a century, and this work tries to add to it. If you are a researcher in nonlinear dynamics, complexity science, or computational systems and you find the framework interesting, the technical paper that follows lays out the mathematics, the stability analysis, the numerical verification, and the falsification criteria in full. The reproducibility archive is public; the equations and predictions are testable on any system that exposes the right operational variables. If you are a reader who came to this essay because the question interested you, and the technical paper that follows is more than you want to read in one sitting, that is fine. The essay you have just read is meant to stand on its own. The paper is here for those who want it. Either way, the conversation is older than this contribution and will continue past it. Schr{\"o}dinger asked what life was. Prigogine asked how order emerged. Kauffman asked where adaptation lived. Friston asked how brains stayed organized. Each of them got something, missed something, and handed the question forward. This paper is an attempt to pass it forward in turn, with one specific contribution and one explicit set of ways it could be wrong. Order holds because something pays the bill. The question is always what the bill is and whether the system can afford it. That question is worth asking precisely. The paper that follows is one attempt at doing so.


\section*{Coda}

The technical paper begins on the next page. It is a short paper --- sixteen pages, with five figures, a falsification table, and two appendices --- and it is written in the spare register that academic physics requires. Readers who have followed the essay this far should find the equations more transparent than they might expect, because the variables and parameters have already been introduced in plain English. The mathematics is doing the work of making the claims precise; the essay has done the work of explaining why the claims matter. If you read only the abstract and the falsification table, you will have the structural argument. If you read the full paper, you will have the proof. If you read neither, you will still, I hope, leave with something --- a clearer sense of how a question about why anything stays together becomes a specific mathematical claim, and what it means to make such a claim in a way that can be checked. Thank you for reading.


\end{document}